\setcounter{topnumber}{2}
\section{Experiments}
\label{sec:experiments}

\subsection{Benchmarks and evaluation setup}
\label{sec:setup}

We evaluate sequential decision-making on Crafter, ALFWorld, and WebShop, and question answering on HotpotQA, WebQuestions, and FinQA \citep{hafner2022crafter,shridhar2021alfworld,yao2022webshop,yang2018hotpotqa,berant2013semantic,chen2021finqa}. Metrics are achievement progress for Crafter, success rate for ALFWorld and WebShop, answer EM/F1 for HotpotQA and WebQuestions, and numeric-answer accuracy for FinQA. \method\ trains LoRA adapters on admitted experience with ETG context \citep{hu2022lora}. Appendix~\ref{app:setup} specifies models, data splits, and training settings. Baseline methods, their source references, and implementation details are provided in Appendix~\ref{app:published}.


\subsection{Task performance}
\label{sec:performance}


Table~\ref{tab:main} reports the progress scores for \method\ on Crafter, ALFWorld, HotpotQA, and WebQuestions using Qwen3-8b. It ranks second on WebShop and FinQA, respectively 0.6 and 1.6 percentage points below the highest listed means. These results cover interactive tasks and textual and numerical reasoning. Appendix~\ref{app:llama31} futher proves \method's robustness using Llama-3.1-8b-Instruct, with results shown in Table~\ref{tab:llama31-8b}.

\vspace{-10pt}
% Requires booktabs, xcolor; existing \method, \bestresult, \secondresult.
% All numeric cells are single-line c columns; no scaling or stacked errors.
\begin{table}[!htbp]
\centering
\footnotesize
\setlength{\tabcolsep}{1pt}
\renewcommand{\arraystretch}{1.15}
\caption{\textbf{Benchmark performance (\%).} Green and yellow indicate the highest and second-highest reported means under the source-specific settings in Appendix~\ref{app:published}. SR: success rate; EM: exact match; Acc.: numeric-answer accuracy.  }
\label{tab:main}

\noindent\makebox[\linewidth][l]{\textbf{(a) Achievement-based exploration}}\par\smallskip
\begin{tabular*}{\linewidth}{@{}l@{\extracolsep{\fill}}ccc@{}}
\toprule
Task / Metric & BALROG & MAGE & \textbf{\method} \\
\midrule
Crafter / Progress / Peak Prog. & $25.5\pm3.2$ & \secondresult{$33.5\pm2.3$} & \bestresult{$\mathbf{34.2\pm1.4}$} \\
\bottomrule
\end{tabular*}
\par\medskip
\noindent\begin{minipage}[t]{0.355\linewidth}
\vspace{0pt}
\textbf{(b) Interaction (SR)}\par\smallskip
\begin{tabular*}{\linewidth}{@{}l@{\extracolsep{\fill}}cc@{}}
\toprule
Method & ALFWorld & WebShop \\
\midrule
ReAct & $57.0\pm3.0$ & $40.0\pm2.4$ \\
Reflexion & $42.7\pm4.0$ & $28.8\pm3.0$ \\
Mem0 & $33.6\pm3.5$ & $2.00\pm0.60$ \\
MemRL & $21.4\pm3.0$ & $9.20\pm1.40$ \\
EvolveR & $73.8\pm2.5$ & $17.6\pm2.0$ \\
GRPO & $77.6\pm2.2$ & $66.1\pm2.0$ \\
SkillRL & \secondresult{$89.9\pm1.8$} & $72.7\pm1.8$ \\
EDGE & $85.9\pm2.0$ & $82.6\pm3.8$ \\
MAGE & $87.3\pm1.8$ & \bestresult{$\mathbf{84.0\pm1.8}$} \\
\method & \bestresult{$\mathbf{90.3\pm1.5}$} & \secondresult{$83.4\pm2.4$} \\
\bottomrule
\end{tabular*}
\end{minipage}\hfill
\begin{minipage}[t]{0.62\linewidth}
\vspace{0pt}
\textbf{(c) Question answering}\par\smallskip
\begin{tabular*}{\linewidth}{@{}l@{\extracolsep{\fill}}cccc@{}}
\toprule
& \multicolumn{2}{c}{HotpotQA} & \multicolumn{2}{c}{WebQuestions} \\
\cmidrule(lr){2-3}\cmidrule(lr){4-5}
Method & EM & F1 & EM & F1 \\
\midrule
CoT & $27.0\pm0.8$ & $39.8\pm0.9$ & $28.8\pm1.0$ & $37.8\pm1.1$ \\
ReAct & $27.6\pm1.2$ & $41.3\pm1.3$ & $27.4\pm1.3$ & $36.3\pm1.4$ \\
EvolveR & $40.3\pm1.0$ & $54.6\pm1.1$ & $36.0\pm1.2$ & $49.8\pm1.3$ \\
SuRe & $38.8\pm0.9$ & $53.5\pm1.0$ & $34.4\pm1.1$ & $48.5\pm1.2$ \\
RPG & $37.9\pm1.1$ & $49.2\pm1.2$ & $34.2\pm1.2$ & $47.3\pm1.3$ \\
EmbQA & \secondresult{$42.0\pm0.8$} & \secondresult{$55.8\pm0.9$} & \secondresult{$38.2\pm1.0$} & \secondresult{$52.1\pm1.1$} \\
RoleRAG & $37.4\pm1.0$ & $49.2\pm1.1$ & $32.7\pm1.2$ & $44.1\pm1.3$ \\
\method & \bestresult{$\mathbf{54.6\pm0.9}$} & \bestresult{$\mathbf{71.5\pm1.0}$} & \bestresult{$\mathbf{40.4\pm1.0}$} & \bestresult{$\mathbf{53.7\pm1.1}$} \\
\bottomrule
\end{tabular*}
\end{minipage}
\par\medskip
\noindent\makebox[\linewidth][l]{\textbf{(d) Financial reasoning (Acc.)}}\par\smallskip
\begin{tabular*}{\linewidth}{@{}l@{\extracolsep{\fill}}ccccc@{}}
\toprule
Task & CoT & MAGE & FinQANet & APOLLO & \textbf{\method} \\
\midrule
FinQA & $57.5\pm2.1$ & $67.5\pm1.7$ & $62.2\pm1.5$ & \bestresult{$\mathbf{70.9\pm1.2}$} & \secondresult{$69.3\pm1.7$} \\
\bottomrule
\end{tabular*}
\end{table}

\vspace{-10pt}

\subsection{Ablation Study}
\label{sec:ablations}


\paragraph{Parameter learning and admission.}
The full method improves over ETG with frozen weights by 3.1--8.4 percentage points across all six benchmarks. SFT without inference-time ETG also exceeds this baseline throughout, indicating that training provides benefits beyond retrieved guidance. Mastery-only admission consistently reduces performance, with the largest decrease on ALFWorld. This supports admitting successful teacher demonstrations before the student reliably performs the associated skill.

\paragraph{Activation and external guidance.}
Removing the activation gate lowers scores on every benchmark. Disabling inference-time ETG after SFT produces larger decreases, particularly on ALFWorld and WebShop. The results support combining learned behavior with continued access to external experience: useful parameter updates do not make retrieval unnecessary for every skill.

\vspace{-10pt}
\begin{table}[!htbp]
\centering\footnotesize
\setlength{\tabcolsep}{2.5pt}
\renewcommand{\arraystretch}{1.05}
\caption{\textbf{Ablation results.} Matched traces and task splits; metrics (\%) are Crafter Progress, ALFWorld/WebShop success, HotpotQA/WebQuestions F1, and FinQA numeric accuracy.}
\label{tab:ablations}
% \resizebox{\linewidth}{!}{%
\begin{tabular}{@{}lrrrrrr@{}}
\toprule
Variant & Crafter & ALFWorld & WebShop & HotpotQA & WebQ & FinQA\\
\midrule
Full \method & \textbf{34.2} & \textbf{90.3} & \textbf{83.4} & \textbf{71.5} & \textbf{53.7} & \textbf{69.3}\\
Mastery-only admission & 30.5 & 85.8 & 79.6 & 68.3 & 51.8 & 66.6\\
ETG with frozen weights & 26.9 & 82.1 & 75.0 & 64.2 & 50.6 & 64.0 \\
Without activation gate & 32.7 & 87.3 & 80.8 & 70.1 & 52.6 & 68.1\\
SFT, no inference ETG & 30.8 & 85.1 & 78.6 & 68.0 & 51.5 & 66.9\\
\bottomrule
\end{tabular}
\end{table}

\vspace{-10pt}

\subsection{Skill verification and memory release}
\label{sec:release-analysis}



\paragraph{Development-set verification.}
Figure~\ref{fig:verification-release} examines a fixed
Qwen3-8B SFT checkpoint with 600 HotpotQA experience
records: 16 original records and 584 added from official
SFT targets.$n$ indicating the number of questions per skill.
Release requires a score strictly above 50\%, together
with the relative-performance and training-evidence checks.

Three of five skills qualify for release.
Complex question parsing retains retrieval because the
candidate scores below Base+ETG.
Multi-answer handling improves over this reference but
fails the absolute performance floor.
These contrasting cases illustrate the complementary
roles of relative and absolute verification.
The resulting decisions archive 331 records and leave
269 retrievable, reducing the retrievable footprint
by 55.2\%.

\begin{figure}[!htbp]
\centering
\includegraphics[width=\linewidth]{figures/verification_release_etg600.pdf}
\caption{\textbf{Skill-level verification and memory release.}}
\label{fig:verification-release}
\end{figure}

\vspace{-5pt}

\paragraph{Held-out evaluation.}
We compare four release policies using the same checkpoint,
initial experience store, and held-out questions:
retain all, immediate release after training, verified
release, and random release at a matched memory budget.
We hold the adapter, prompts, retrieval configuration,
and decoding settings fixed and freeze release decisions
before test evaluation.
This comparison isolates release-policy effects from
parameter learning and adapter activation.
Appendix~\ref{app:release-controls} provides the detailed
protocol and uncertainty estimation.


Table~\ref{tab:release-controls} shows that verified release
remains close to retain-all performance while making
fewer than half of the original records retrievable.
Its paired F1 difference is $-0.2$ points, with a 95\%
confidence interval of $[-0.6,\,0.2]$.
This interval is consistent with a small performance
change in either direction, rather than establishing
that release causes no loss.

Immediate release leaves only 13 retrievable records
but lowers F1 from 71.1 to 66.8, indicating that training
participation alone is insufficient to justify aggressive
release in this setting.
At the same retrievable record count, verified release
outperforms random release by 3.3 F1 points.
EM follows the same ordering.
These results support verification as a means of reducing
retrieval dependence while limiting the observed
performance cost.
All released records remain archived, so the reduction
concerns retrieval availability rather than total storage.

% Requires: booktabs
\begin{table}[t]
\centering
\footnotesize
\setlength{\tabcolsep}{4pt}
\renewcommand{\arraystretch}{1.1}

\caption{\textbf{Held-out performance under alternative
memory-release policies.}
All conditions use the same HotpotQA checkpoint,
initial ETG store, and test questions.
EM and F1 are percentages.
$\Delta$F1 is the paired difference from retaining all
records with a 95\% confidence
interval.}
\label{tab:release-controls}

\begin{tabular*}{\linewidth}{
@{}l@{\extracolsep{\fill}}ccccc@{}}
\toprule
& \multicolumn{3}{c}{Held-out performance}
& \multicolumn{2}{c}{ETG records}\\
\cmidrule(lr){2-4}
\cmidrule(lr){5-6}
Release policy
& EM
& F1
& $\Delta$F1 [95\% CI]
& Retrievable
& Archived\\
\midrule
Retain all
& 54.4
& 71.1
& Reference
& 600
& 0\\

Immediate release
& 50.9
& 66.8
& $-4.3\;[-4.6,\,-4.0]$
& 13
& 587\\

Random release
& 51.6
& 67.6
& $-3.5\;[-3.9,\,-3.1]$
& 269
& 331\\

Verified release
& 54.0
& 70.9
& $-0.2\;[-0.6,\,0.2]$
& 269
& 331\\
\bottomrule
\end{tabular*}
\end{table}
\vspace{-5pt}

\subsection{Latency and serving cost}
\label{sec:efficiency}

\vspace{-5pt}
\begin{table}[!htbp]
\centering\footnotesize\setlength{\tabcolsep}{4pt}
\caption{\textbf{Latency and serving work.} Means per task on paired ALFWorld and WebQuestions using Qwen3-8B. \colorbox{enggreen}{\textbf{Green}} and \colorbox{engyellow}{yellow}: highest and second-highest values. Lower is better for time, calls, and token counts; higher is better for SR/F1. Input/output tokens are in thousands. Reflexion includes up to three attempts and reflection calls. Dashes denote pending supplementary measurements.}
\label{tab:latency}
\begin{tabular}{@{}llrrrrr@{}}
\toprule
Benchmark & Method & Time (s) & Calls & Input (k) & Output (k) & SR / F1 (\%; $\Delta$ pp)\\
\midrule
ALFWorld & ReAct-style & \secondresult{50.27} & 36.50 & 73.64 & 1.95 & 57.0(-33.3) \\
 & Reflexion & 167.86 & 104.33 & 252.86 & 6.92 & 42.7(-47.6)\\
& Base + ETG
& \bestresult{\textbf{43.76}}
& 29.50
& 81.24
& 1.66
& $82.1\,(-8.2)$\\
& SFT, no inference ETG
& 60.53
& \bestresult{\textbf{22.58}}
& \bestresult{\textbf{35.76}}
& \bestresult{\textbf{1.30}}
& \secondresult{$85.1\,(-5.2)$}\\
 & \textbf{\method}
& 65.91
& \secondresult{25.50}
& \secondresult{46.50}
& \secondresult{1.51}
& \bestresult{\textbf{90.3}\,$(+0.0)$}\\
\midrule
WebQuestions & CoT & 1.56 & \bestresult{\textbf{1.00}} & 1.74 & 0.07 & 37.8(-15.9) \\
 & Self-critique (3 calls) & 8.83 & \secondresult{3.00} & 5.61 & 0.43 & 40.6(-13.1)\\
 & Base + ETG & 1.67 & \bestresult{\textbf{1.00}} & 1.78 & 0.07 & 50.6(-3.1)\\
 & SFT, no inference ETG & \bestresult{\textbf{0.86}}
& \bestresult{\textbf{1.00}}
& \bestresult{\textbf{1.74}}
& \bestresult{\textbf{0.02}}
& \secondresult{$51.5\,(-2.2)$}\\
 & \method & \secondresult{1.26} & \bestresult{\textbf{1.00}} & \secondresult{1.75} & \secondresult{0.04} & \bestresult{\textbf{53.7}\,$(+0.0)$} \\
\bottomrule
\end{tabular}
\end{table}



On ALFWorld, \method\ uses 42.8\% fewer input tokens and 13.6\% fewer calls than Base+ETG, but takes longer overall. SFT without inference-time ETG uses still fewer resources, while \method\ improves success by 5.2 percentage points. Thus, the full controller trades additional inference work for higher task performance. On WebQuestions, \method\ is 24.6\% faster than Base+ETG and improves F1 by 3.1 points. SFT without inference-time ETG is the fastest configuration, but scores 2.2 F1 points below \method. Together, these results show that \method\ achieves the strongest task performance among the measured configurations, while the latency trade-off depends on the benchmark and comparison method.
